%As a result, the concerns that define MP-VRDU receive no systematic treatment: how architectures manage cross-page evidence under context-window constraints; how text, visual, and structural information trade page coverage against reading fidelity; how inference-time retrieval, reasoning, and agentic orchestration coordinate; how training strategies adapt to document-scale supervision; and how LLM-synthesised data shapes both training and evaluation integrity.
% Two trends emerge, distinguished by where the document resides at inference, and each divides into two families. \textbf{In-context encoding} holds the whole document within a single model's forward pass, covering page-to-document encoding models (\S\ref{sec:arch-p2d}) and backbone-centric adaptation models (\S\ref{sec:arch-mllm}). \textbf{Selective evidence pipelines} instead pass only a compact evidence subset to a generator, covering retrieval-augmented (\S\ref{sec:arch-ra}) and agentic pipelines (\S\ref{sec:arch-agent}). %The two trace a broadly chronological shift from encoding all pages jointly toward selectively retrieving and iteratively inspecting evidence.
% Page-to-document encoding models introduce explicit architectural primitives for cross-page exchange, producing document-level outputs in a single forward pass~\cite{tito2023hierarchical,blau2024gram}. Those frameworks carry the multi-page mechanism in the architecture itself through primitives such as page summary tokens, document-global attention, recurrent memory, or sparse cross-page attention~\cite{dong2024multi,borchmann2025arctic}. The resulting inductive bias for multi-page structure suits documents with regular templates such as forms or reports, but the design pays an intrinsic cost since all pages still pass through a shared encoder, compute scales with document length, and per-page compression discards token-level detail. As pretrained MLLMs with widened context windows now absorb cross-page evidence implicitly, the practical motivation for dedicated cross-page architectures has weakened.
% \textbf{Page-to-Document Encoding Models} encode each page through a shared document encoder including page level OCR text, layout, and visual features, then connect page representations through a document-level exchange mechanism such as summary tokens \cite{}, document-global tokens \cite{}, recurrent memory \cite{}, or sparse cross-page attention \cite{}. Cross-page interaction occur within the framework in a single forward pass, while page identity, local layout, and reading order remain explicit, providing a strong inductive bias for multi-page structure. Explicit cross-page modelling makes these models well suited to tasks where relevant evidence appears across pages within a coherent document structure, such as reports and business documents \cite{tito2023hierarchical, borchmann2025arctic}. Its limitations follow from the same design, as document-level representations may lose-fine grained evidence when page information is compressed or weakly connected \cite{blau2024gram, dong2024multi}. Moreover, computation grows with document length as all pages typically need to be encoded, making very long documents difficult to scale.
% MLLM-centric adaptation models directly process the document inside a pre-trained MLLM whose context window absorbs cross-page information directly. Adaptation strategies range from training only lightweight projectors~\cite{tanaka2024instructdoc} or visual-token compression modules~\cite{hu2025mplug2}, to fine-tuning the backbone itself through continued pre-training~\cite{duan2025docopilot}, instruction tuning, or reinforcement learning~\cite{xiong2026docr1}. This design benefits directly from the MLLM's instruction-following and reasoning capabilities. However, the backbone's context length now caps total document size, fine-grained signals depend on the fidelity of token compression~\cite{jia2024leopard}, and document structure such as section hierarchy and figure-caption linkage must be recovered implicitly~\cite{hannan2025docslm}.
%, limits that motivate moving evidence selection out of the backbone altogether.
% This family adapts a pretrained MLLM backbone to multi-page input by training projectors that map per-modality features into the backbone's token space \cite{tanaka2024instructdoc}, adaptors that enable parameter-efficient fine-tuning for document understanding \cite{zhu2025laytokenllm}, and compression modules that contain visual token growth on long inputs \cite{hu2025mplug2}. Some works further unfreeze the backbone for continued pretraining or instruction tuning on document corpora, trading parameter efficiency for stronger document-specific grounding \cite{duan2025docopilot}. These models inherit the instruction-following and multimodal reasoning of the foundation model and remain parameter-efficient to train, fitting tasks that combine document reading with open-ended reasoning and moderate-length documents where long-context tuning massively improves MLLM performance on multi-page documents \cite{hu2024mplug, yu2024texthawk2}. However, as the backbone must absorb document-scale evidence internally, performance degrades beyond moderate-length documents, leaving quality bound to context length and to how faithfully fine-grained text and visual signal survive token compression \cite{jia2024leopard}. Additionally, reasoning over visual elements depends entirely on the backbone's pretrained visual capabilities, and as most models in this family are OCR-free, reliable text recovery requires high-resolution input that aggressive downsampling readily compromises. Moreover, document structure such as section hierarchy and figure-caption linkage must likewise be recovered implicitly from pixels rather than supplied through layout tokens or a structural encoder \cite{hannan2025docslm, duan2025docopilot}.
% Retrieval-augmented pipelines first select relevant evidence units, such as chunks, pages, or graph nodes, and then condition the generator on this compact subset.~\cite{cho2024m3docrag,tanaka2025vdocrag}. This separation between evidence selection and answer generation allows them to scale to longer multi-page inputs than MLLM-centric and page-to-document encoding models. The two stages may run as a single forward pass or as a fixed arrangement of specialised readers~\cite{han2025mdocagent}, modalities-specific branches~\cite{suri2025visdom}, or synthesis components~\cite{duan2025m2rag}, but in either case, the control flow is set at design time. This fixed structure means omitted or poorly ranked evidence is difficult to recovered during generation, aggregation over multiple weakly relevant pages remains challenging~\cite{wu2025molorag}, and upstream errors in chunking, OCR, or modality alignment can directly affect the final answer.
% If want to go with agentic, could name it "Adaptive-Agent Pipelines"
% May need to reduce length if required
% Adaptive-trajectory pipelines place the trajectory under model control, as one more more LLM controllers plan the sequence of evidence-gathering actions at infernece time and select each action from the outcomes of earlier ones rather than executing a plan fixed in advance.~\cite{zheng2026docvstar,sun2025docagent}. A controller action may surface as an explicit call to an external retriever, parser, or OCR engine~\cite{wang2025vidorag,yu2025mact}, or a structured response the framework interprets as a re-retrieval or as navigation over a pre-built document representation~\cite{jain2025simpledoc,fu2026dmap}. Roles such as inspection, synthesis, and adjudication may again be distributed across multiple backbone LLMs, but the defining property is that the trajectory adapts to question difficulty, as such a multi-component system with a predetermined execution order remains in the retrieval-augmented category even when its components are described as ``agents''. Model-controlled planning lets compute scale with question difficulty rather than document length and exposes the full action sequence, not just the final evidence, as an inspectable audit trail. However, the variable-length trajectory makes inference cost difficult to bound in advance, controllers are sensitive to prompt phrasing and may commit early to wrong evidence and fail silently, stochastic generation complicates reproducibility, and the system still inherits the recall ceiling of every retriever, parser, or scoring module it consults.
% \subsection{Agentic Pipelines}
% Agentic pipelines replace the fixed retrieval pipeline with iterative evidence acquisition: one or more LLM controllers issue successive tool calls, for retrieval, navigation, cropping, or OCR, and each subsequent decision is conditioned on the outcomes of earlier observations~\cite{zheng2026docvstar,sun2025docagent}. Roles such as inspection, synthesis, and adjudication may be distributed across heterogeneous backbones~\cite{yu2025mact,wang2025vidorag}, but the defining property of an agentic system is that its reasoning trajectory adapts to question difficulty rather than following a fixed control flow; multi-component systems with predetermined execution order, even when described as ``agents'', remain within the retriever-generator family. By lifting the single-retrieval-step bottleneck of retriever-generator systems, this design allows compute to scale with question difficulty rather than document length and exposes the trajectory as an inspectable audit trail. However, the variable-length action sequence makes inference cost difficult to bound before execution, controllers remain sensitive to prompt phrasing and may commit early to wrong evidence with silent failure, reproducibility is complicated by stochastic generation, and the system inherits the recall ceiling of every retriever, parser, or OCR module it calls.
% Recent surveys \cite{plaat2025agentic, wang2024survey} characterise agentic LLM systems through their capacity to reason, act, and interact with an external environment beyond a single forward pass. Following this framing, we categorise MP-VRDU systems that perform multi-pass evidence acquisition or distributes work across specialised agents as agentic, while reserving the retriever-generator family for single-pass models. 
% multi-page settings introduce distinct challenges for text, vision, and document structure: text may exceed the backbone context window and appear sparsely across pages; visual processing becomes expensive as both page number and resolution increase; and document-level structures such as section hierarchies, multi-page tables, and cross-references cannot be captured by single-page layout signals alone. We therefore discuss text, visual, and structural modelling in turn (\S\ref{sec:modality-text}, \S\ref{sec:modality-visual}, \S\ref{sec:modality-structure}), focusing on the issues specific to MP-VRDU rather than those already covered in single-page surveys~\cite{ding2025survey}. Modality fusion is treated as a complementary cross-cutting issue in Appendix~\ref{app:fusion}.
% Text in MP-VRDU presents three bottlenecks. \emph{Volume} across  pages normally exceeds what the MLLM can process in full; \emph{continuity} breaks when paragraphs, tables, and section context span page boundaries; and \emph{heterogeneity} in text density varies sharply across pages within a single document. 
% \textbf{Volume} is addressed either by compressing all pages into a fixed budget through page summaries, memory tokens, or long-context encoders~\cite{borchmann2025arctic, dong2024multi}, or by indexing chunks and passing only selected text to the reader~\cite{xie2024wukong, zhang2024cream}. The first fixes compression before the question is seen; the second makes evidence missed at retrieval unrecoverable. \textbf{Continuity} is preserved by recurrent state propagation that 
% carries page context forward~\cite{dong2024multi}, structure-aware chunking that groups text under recovered section headings~\cite{shin2025multidocfusion}, and adjacent-page expansion at retrieval time~\cite{qian2026accurate}, each restoring local cross-page context at the cost of broader query relevance. \textbf{Heterogeneity} receives the least attention: most systems apply fixed token budgets and fixed top-$k$ uniformly, over-compressing dense pages while wasting capacity on sparse ones, with adaptive selection~\cite{li2025avir} only beginning to address per-query variability.
% Visual evidence preserves signals that text extraction loses, but in MP-VRDU its cost structure introduces three bottlenecks. \emph{Resolution} cannot be maintained at full document length, since token cost scales with page count multiplied by per-page resolution; \emph{multi-signal competition} forces text-bearing pixels, non-textual evidence (charts, figures, handwriting), and layout cues to share the same budget on each page; and \emph{heterogeneity} in visual complexity varies sharply across pages within a single document.
% \textbf{Resolution} forces a choice between compression and selective high-fidelity reading. Token-reduction modules and adaptive sub-image budgets~\cite{hu2025mplug2, yu2024texthawk2, jia2024leopard} keep every page available but leave fine-grained content vulnerable to downsampling, while inspecting only a retrieved subset at high resolution~\cite{zhu2025doclens, yan2025docseeker} preserves detail at the cost of inheriting the selector's recall ceiling. \textbf{Multi-signal competition} arises because the same visual budget need support text, table and chart understanding. Visual RAG systems encode rendered pages directly~\cite{cho2024m3docrag, tanaka2025vdocrag} so all signals share one representation but the page becomes a single evidence unit, whereas element-level cropping of tables, figures, and charts~\cite{zhu2025doclens} preserves fine structure at the cost of detector dependence and separation from surrounding explanatory text. \textbf{Heterogeneity} across pages receives the least attention: most systems apply uniform resolution and sub-image budgets regardless of how visually dense a page is, and evidence-guided allocation~\cite{yan2025docseeker, jia2024leopard} only begins to address per-page and per-query variability.
% Structure turns isolated page observations into document evidence, and in MP-VRDU the difficult structure is not token geometry within one page but relationships whose scope crosses pages. \textbf{Hierarchy and reading order} is the foundational problem: sections, lists, columns, and captions routinely span page boundaries. Explicit parsers expose this as outlines or hierarchical indices~\cite{shin2025multidocfusion, sun2025docagent}, while OCR-free backbones attempt implicit recovery from pixels~\cite{yu2024texthawk2}; the unresolved gap is robustness, since parser errors propagate downstream and implicit recovery depends on high-resolution input that compression may have already removed. \textbf{Relations and references} demand following typed links rather than ranking similar pages. Graph and map-based methods make these relations first-class objects~\cite{wang2024knowledge, wu2025molorag, fu2026dmap}, improving navigability at the cost of edge quality bounded by extraction heuristics or LLM judgement. \textbf{Continuation} extends the same problem to multi-page tables and section-spanning arguments, where the relevant unit is larger than a page but smaller than the document. Hierarchical and cross-page retrieval preserve parent context~\cite{gong2025mhier, qian2026accurate}, but the boundary between useful context and irrelevant neighbours remains hard to set.
% Sparse lexical signals such as BM-25 serve as baselines or auxiliary channels, while dense retrieval in MP-VRDU typically uses vision-language late-interaction encoders (ColPali, ColBERT-v2, ColQwen) that align token-level query vectors with patch-level page embeddings, preserving fine-grained alignment between query terms and page regions. M3DocRAG~\cite{cho2024m3docrag} retrieves at chunk and page granularity with this signal alone, while AVIR~\cite{li2025avir} adds an adaptive selector for per-query page counts. Others add a second reranking stage: CREAM~\cite{zhang2024cream} pairs BGE retrieval with iterative LLM reranking, SimpleDoc~\cite{jain2025simpledoc} combines ColPali embeddings with LLM summary reranking under persistent working memory, and MDocAgent~\cite{han2025mdocagent} aggregates parallel text and image retrieval tracks downstream. Independent scoring handles paraphrase and cross-modal queries well, but cannot follow cross-page references or multi-hop dependencies that no single passage carries.
% Relation-aware methods retrieve along explicit document structure, rather than scoring passages independently against the query. Tree-based variants such as MultiDocFusion~\cite{shin2025multidocfusion} and DocAgent~\cite{sun2025docagent} extract section hierarchies, outlines and group passages along semantic boundaries so retrieved chunks preserve parent context, while MHier-RAG~\cite{gong2025mhier} extends this by routing queries through dual in-page and cross-page indices. Graph-based variants treat headings, entities, and cross-references as graph nodes with typed edges, allowing controllers to traverse from one passage to related ones rather than relying on similarity alone: KGP~\cite{wang2024knowledge} uses an LLM as a traverser over a passage-and-structure graph, while MoLoRAG~\cite{wu2025molorag} and MLDocRAG~\cite{zhang2026mldocrag} extend this with logical relevance scoring or query-conditioned graph construction. Structure-based methods recover references that similarity scoring misses, at the cost of dependence on parser quality and traversal cost that grows with the hop count.
% \paragraph{Sparse Hierarchy-Aware Search.}
% Many models parse the input document into an explicit textual structure and perform sparse retrieval or traversal through term or structure anchored matching. MultiDocFusion~\cite{shin2025multidocfusion} uses a hierarchy parser to recover document, section, and subsection trees from scanned PDFs and groups passages along DFS-ordered Markdown headers, so retrieved chunks preserve their parent context (parser tuned). DocAgent~\cite{sun2025docagent} instead extracts a tree-formatted XML outline from headings, sections, and captions and exposes it through section-level search and section-content fetch tools that navigate the parsed hierarchy directly. MHier-RAG~\cite{gong2025mhier} extends the structure to a dual hierarchy of in-page flattened chunks and cross-page topological chunks, with LLM-based re-ranking across both indices, addressing the cross-page fragmentation that single-tree designs leave unresolved. However, sparse signals match on lexical or structural overlap and miss semantic connections and paraphrased queries, resulting in degraded performance on questions whose phrasing diverges from the document's surface form.
% \paragraph{Dense Embedding Search.}
% Dense embedding retrieval addresses this by encoding pages or chunks into learned representations, so queries match on semantic similarity rather than lexical or structural overlap. Many models \cite{wu2025docreact, wu2025molorag, han2025mdocagent, jain2025simpledoc, wang2025vidorag} adopt vision-language late-interaction encoders such as ColPali, ColBERT, and ColQwen as their dense retriever. These encoders embed the query as a sequence of contextual token vectors and each page as a bag of patch embeddings produced by a vision-language transformer over the rendered page image, then score a query-page pair by taking the maximum cosine similarity of each query token over all page patches and summing the result. MaxSim aggregation preserves fine-grained alignment between specific query terms and specific page regions, which is critical in MP-VRDU as a question token may align with a single figure caption, table cell, or paragraph buried in a long document. M3DocRAG~\cite{cho2024m3docrag} instantiates this pattern by retrieving at chunk and page granularity with the rest of the pipeline frozen. AVIR~\cite{li2025avir} adds an adaptive page selector that examines the per-query score distribution and either thresholds at a fixed cutoff for short documents or applies K-means clustering with a top-K bound on longer ones, so the number of pages handed to the reasoner scales with the score distribution rather than being fixed (lightweight head trained). Such methods handle paraphrase and cross-modal queries well, yet remain structure-blind on their own, since a section that should be retrieved as a unit can be split across embedding chunks and cross-page references are not represented as first-class objects.
% \paragraph{Joint Multi-Signal Search.}
% The complementary strengths of sparse and dense signals can be combined at the retrieval stage rather than committing to a single channel. CREAM pairs BGE coarse retrieval with iterative LLM-based re-ranking across overlapping chunk groups, so dense semantic candidates are filtered by an LLM judgement that draws on the full passage text \cite{zhang2024cream}. SimpleDoc combines ColPali page embeddings with LLM-generated summary re-ranking under a working memory that persists across rounds, so successive iterations can refine the joint score \cite{jain2025simpledoc}. MDocAgent runs parallel text and image retrieval tracks whose candidates are aggregated by a downstream agent, joining textual and visual evidence rather than fusing modalities at encoding time \cite{han2025mdocagent}. Such combinations improve recall by spanning what any single signal misses, but still score each candidate passage in isolation, so retrieval cannot follow logical links, cross-references, or multi-hop dependencies that span passages.
% \paragraph{Graph-Based Traversal}
% Graph traversal addresses this by indexing the document as a graph and navigating it through edges, resulting in strong inter-passage. KGP~\cite{wang2024knowledge} constructs a knowledge graph over passages and document-structure nodes such as pages, paragraphs, and tables, with edges encoding lexical similarity and structural belonging, then uses an instruction-tuned LLM as a traverser that predicts the next evidence node from the visited path. MoLoRAG~\cite{wu2025molorag} instead builds a page graph in which edge candidates are filtered by embedding similarity, then scores logical relevance with a frozen Qwen2.5-VL prompt and combines that with semantic similarity for multi-hop traversal, so retrieval can follow logical dependencies that surface similarity misses. MLDocRAG~\cite{zhang2026mldocrag} constructs the graph per query, anchoring chunk nodes to LVLM-generated query nodes and expanding through KNN edges between queries, so the structure adapts to the question rather than being fixed at index time. Graph methods reach references that scoring methods cannot, but results in construction overhead, edge-quality dependence, and traversal cost that scales with hop count.
%These patterns are independent of architectural family: they can wrap a single retrieval-augmented forward pass or operate inside one step of an agentic loop, in which case the loop and the within-step pattern compound their effects on robustness and cost.
% \paragraph{Chain-of-Thought and Self-Improvement.}
% Chain-of-thought prompting elicits intermediate reasoning steps from a frozen LLM or VLM that integrate retrieved text, tables, and figures and combine facts across pages before producing the final answer. Following this definition by \cite{wei2022chain}, the entire reasoning chain runs over a fixed evidence buffer assembled by a prior retrieval step or over the entire compressed document, so the reasoner cannot acquire new or finer grained evidence once decoding begins. MHier-RAG~\cite{gong2025mhier} applies a structured CoT prompt to its hierarchical retrieval output, although a single reasoning chain remains vulnerable to compounding errors on dense layouts. SimpleDoc~\cite{jain2025simpledoc} interleaves self-refinement and self-verification, emitting a refined query when evidence is judged insufficient, although the loop terminates on the backbone's own judgement and inherits its blind spots. DocSLM~\cite{hannan2025docslm} (trained, see \S4) instead applies document-level self-consistency, streaming page segments and selecting the lowest-uncertainty answer from a learned uncertainty head, trading compute for robustness against single-pass errors. DocR1~\cite{xiong2026docr1} (trained, see \S4) internalises the CoT trajectory through reinforcement learning, producing structured \texttt{<think>}, \texttt{<evidence\_page>}, and \texttt{<answer>} outputs that interleave reasoning with explicit evidence localisation. Across these variants, reasoning operates over a fixed evidence buffer, so questions whose evidence sits outside the retrieved set remain unreachable regardless of how the trace is sampled or verified.
% \paragraph{ReAct and Tool-Augmented Reasoning.}
% ReAct~\cite{yao2022react} interleaves reasoning and acting, alternating intermediate reasoning steps with tool calls whose observations feed the next step. In MP-VRDU the tools retrieve, fetch, or inspect document pages, letting the reasoner gather evidence during inference rather than work from a closed buffer. Doc-React~\cite{wu2025docreact} instantiates flat-retrieval ReAct with InfoNCE-guided sub-query refinement and entropy-based termination, although each iteration still selects from a single retriever's ranked list and recovery from a poor encoder match relies on subsequent query rewrites. Outline-driven variants combine this loop with the structured navigation of \S5.1. DocAgent~\cite{sun2025docagent} traverses its tree-formatted XML outline through section-level search and fetch tools chosen turn-by-turn by an actor agent, and DocDancer~\cite{zhang2026docdancer} (trajectory-trained) extends this with explicit intent modelling so action choices are conditioned on the agent's evolving goal, although both inherit the parser's outline quality and degrade on documents without exploitable structure. Doc-V$^*$~\cite{zheng2026docvstar} (RL-trained) couples ReAct with a thumbnail-level document map, alternating semantic search with direct page indexing under a working memory of summaries. ReAct loops sharpen evidence localisation through dynamic action selection, yet pay in tool latency that grows with rounds, prompt sensitivity at each step, and a fixed-resolution reading of any fetched page that leaves fine-grained content exposed to attention dilution.
%Tool latency, compounding error from early wrong actions, and unbounded trajectory length remain shared limitations.
% \paragraph{Modality-Specialised Collaboration.}
% Modality-specialised orchestration assigns separate agents to text and visual evidence and adds a synthesis agent that reconciles their outputs, addressing the observation that text-only and image-only readers each miss evidence the other captures. The pattern preserves single-modality reading fidelity but depends on the synthesis step to resolve disagreements without access to the original document. MDocAgent~\cite{han2025mdocagent} runs parallel text and image RAG with a critical agent that extracts cross-modal cues for downstream text and image specialists, while M$^2$RAG~\cite{duan2025m2rag} fuses a text filter and visual extractor through a modality-fusion agent that arbitrates between modalities.
% \paragraph{Role-Specialised Procedural Decomposition.}
% Role-specialised orchestration decomposes the procedural stages of document reasoning, such as planning, execution, verification, and synthesis, into dedicated agents, with explicit routing between roles when a downstream agent detects an error. Separating judgement from generation reduces the cognitive blind spots of self-correction, but introduces inter-agent compatibility costs and a recall ceiling bounded by upstream stages. MACT~\cite{yu2025mact} instantiates planning, execution, judgement, and answer agents with the judgement agent rerouting flawed plans or executions to their originator, while DocAgent~\cite{sun2025docagent} pairs an actor that navigates a parsed outline through tool calls with a reviewer that validates the initial response against cross-modal evidence.
% Note on training: pretraining/continued-pretraining is rarely done in MP-VRDU, mostly SFT and IT for training trajectory, reasoninig, orchestration etc.
%While \S\ref{sec:inference-time} held the backbone fixed, many MP-VRDU systems extend these pipelines by training specific components. This section instead organises MP-VRDU training by the inference-time component being trained, namely the retriever (\S\ref{sec:train-retrieval}), the reasoner (\S\ref{sec:train-reasoning}), and the agentic orchestration policy (\S\ref{sec:train-agentic}). Two cross-cutting techniques (\S\ref{sec:train-crosscutting})
% Training strategies in MP-VRDU are closely tied to how each architecture controls the cost and complexity of multi-page document understanding. While the previous section describe how models represent documents and combine text, visual, and layout evidence, this section focuses on how these capabilities are learned.
% Prompted reasoning leaves quality sensitive to prompt design and base-model behaviour. Reward-driven variants train a reasoning policy with explicit evidence supervision: DocR1~\cite{xiong2026docr1} fine-tunes Qwen2.5-VL with Evidence-Page-Guided GRPO, whose reward combines format compliance, answer accuracy, and a page-citation indicator, producing structured \texttt{<think>}, \texttt{<evidence\_page>}, and \texttt{<answer>} outputs. Trace-imitation variants instil a reasoning chain through SFT and preference optimisation: CoR~\cite{guo2026end} bakes in a four-stage plan-execute-integrate-verify chain through LoRA SFT on a mixed corpus, full-parameter SFT on the 26K CoR-Dataset and Mask-AR, then DPO on 5K preference pairs. Uncertainty-driven variants train an aggregation head rather than a trajectory: DocSLM~\cite{hannan2025docslm} learns an entropy-based calibrator that selects the lowest-uncertainty answer across streaming segments, providing a self-consistency signal under tight parameter budgets.
% Prompt-only orchestration is unreliable on smaller open-weight backbones that cannot follow ReAct protocols. Behavioural-cloning variants distil agent trajectories from a stronger model: DocDancer~\cite{zhang2026docdancer} fine-tunes a 4B or 30B Qwen3 agent on synthesised exploration-then-synthesis traces that pair ReAct \texttt{Search} and \texttt{Read} actions with intent-conditioned tool selection. Reinforcement-learned variants reward correct tool use and trajectory shape: Doc-V*~\cite{zheng2026docvstar} couples SFT on coarse-to-fine ReAct traces with GRPO over a two-action space of \texttt{retrieval\_page} and \texttt{fetch\_page} calls plus entropy-based stopping, while MACT~\cite{yu2025mact} trains a multi-agent policy with agent-wise adaptive test-time scaling that allocates sample budgets per role of planner, executor, judge, and answerer by per-question difficulty.
% Copied from Lit Review haha
% \subsection{Parameter Efficient Tuning}
% \label{sec:train-peft}
% Full backbone fine-tuning is prohibitive at long-document context lengths and erases pretrained capabilities, so the dominant MP-VRDU pattern freezes the backbone and trains lightweight modules such as LoRA matrices, projection heads, or task-specific adapters. LayTokenLLM~\cite{zhu2025laytokenllm} freezes the LLM and adds a layout tokenisation module plus LoRA, encoding spatial layout through compressed positional tokens with negligible context overhead, while CREAM~\cite{zhang2024cream} applies the LLaMA-Adapter-V2 recipe to a LLaMA2-7B answerer, training only LoRA matrices, attention pooling, and bias and normalisation parameters above a frozen Pix2Struct encoder. The same pattern recurs in DFVC's frozen-ColPali adapter (\S\ref{sec:train-retrieval}) and DocDancer's LoRA trajectory tuning (\S\ref{sec:train-agentic}), confirming that PEFT versus full fine-tuning is more often a compute-budget decision than a target-component one.
% \subsection{Curriculum / Multi-Stage Training}
% \label{sec:train-curriculum}
% Several systems converge on a multi-stage curriculum that starts on single-page or local objectives before introducing document-scale supervision, exploiting the fact that text reading and layout understanding transfer from single-page data while cross-page integration must be learned separately. mPLUG-DocOwl2~\cite{hu2025mplug2} uses three stages of single-image structure learning on DocStruct4M, MP-DocStruct1M continue-pretraining, and multi-task fine-tuning, while DocSLM~\cite{hannan2025docslm} adds two more stages for a five-stage schedule covering image-OCR alignment, image- and document-level compression, instruction tuning, and streaming-abstention calibration.
% Datasets in MP-VRDU span three roles. Pretraining historically relied on single-page archives such as IIT-CDIP \cite{lewis2006iitcdip} and OCR-IDL \cite{biten2022ocridl}, with genuinely multi-page corpora only recently emerging through DocStruct4M's mixed single- and multi-page pretraining \cite{hu2024mplug} and the purely multi-page MP-DocStruct1M, Doc-750K, and scientific PaperPDF \cite{hu2025mplug2, duan2025docopilot, xie2024wukong}. Fine-tuning has shifted toward bespoke synthetic generation per model, with strong LVLMs writing question-answer or trajectory data over scraped PDFs since per-page annotation does not scale to long documents \cite{tanaka2024instructdoc, jain2025simpledoc, xiong2026docr1, zhang2026docdancer}. Benchmarks have evolved from mid-length extractive evaluation on MP-DocVQA, DUDE, and SlideVQA \cite{tito2023hierarchical, van2023dude, tanaka2023slidevqa} to long-context multimodal evaluation on MMLongBench-Doc, LongDocURL, and DocBench \cite{ma2024mmlongbench, deng2025longdocurl, zou2024docbench}, with retrieval and page-localisation increasingly scored alongside answer accuracy \cite{hui2024uda, xiong2026docr1, zhu2025doclens}. Five benchmarks (MP-DocVQA, MMLongBench-Doc, DUDE, DocVQA, and InfographicVQA) appear in roughly half or more of the surveyed papers and constitute the de-facto axes for cross-model comparison. A detailed treatment with per-dataset statistics is provided in Appendix~\ref{dataset_detailed}, the dataset catalogue in Table~\ref{tab:dataset_survey}, and the per-paper usage matrix in Table~\ref{tab:datasets_used}.
% MP-VRDU systems must reconcile high-resolution per-page reading with document-scale evidence aggregation, an objective that exceeds any MLLM's compute and context budget~\cite{ma2024mmlongbench}. Every architectural family takes a lossy compromise: compression silently drops fine-grained content and is typically fixed at index time so evidence is discarded before the question is seen; retrieval bounds answer accuracy from above while broadening the retrieved set introduces attention dilution; and iterative retrieval lifts this ceiling only partially, inheriting the recall of every tool it calls. Cross-page structure such as section hierarchies and inter-page references compounds this further, as none of the three current routes (\S\ref{sec:modality-structure}) scales cleanly. Future work should explore query-conditional and modality-aware compression that retains evidence based on a question's evidential demands, retrieval that follows cross-references and logical relations rather than scoring passages in isolation, and learned cross-page structural representations that survive token compression.
% Beyond contamination, two further evaluation gaps persist. Answer-only accuracy masks retrieval failure and rewards correct answers reached for the wrong reason; newer benchmarks score retrieval and page-localisation alongside answers~\cite{hui2024uda, xiong2026docr1, zhu2025doclens}, but the practice is not yet standardised, and agentic pipelines produce stochastic trajectories that complicate reproducibility. Task and language scope also remain narrow: surveyed systems are overwhelmingly DocVQA-centric and English-centric, while real-world MP-VRDU needs span summarisation, cross-document reasoning, and structured extraction across diverse non-English layouts and writing systems where reading order, script direction, and document conventions differ substantially. Future work should standardise trajectory-level evaluation paired with retrieval and grounding metrics, report stochasticity bounds for agentic systems, and broaden benchmark coverage across task types, document domains, and languages, particularly low-resource and non-Latin-script settings that current pretraining corpora barely touch.
% \section{Conclusion}
% This survey reframes multi-page visually rich document understanding around cross-page evidence management, organising architectures, modalities, inference-time adaptation, and training strategies by how each handles evidence that no single forward pass can contain. A consistent pattern emerges across families: progress comes from moving evidence handling out of the backbone — through compression, retrieval, or iterative orchestration — but each shift relocates the bottleneck rather than removing it. Closing the remaining gaps in efficiency, faithfulness, evidence pipelines, supervision integrity, and evaluation scope is the open frontier for the next generation of MP-VRDU systems.
% \section*{Acknowledgements}
% Bibliography entries for the entire Anthology, followed by custom entries
%\bibliography{anthology,custom}
% Custom bibliography entries only
% Not that great, could omit entirely.
% \paragraph{Staged Fusion.}
% Local alignment is combined with selective document-level reasoning. Modalities are fused within pages or regions, then compressed, retrieved, or stored selectively before a final cross-page fusion step. Some hierarchical and recurrent models approximate the pattern by producing page-level summaries or memories prior to document-level aggregation \cite{tito2023hierarchical, dong2024multi}, and compact long-document models reduce local multimodal evidence into fixed-size page representations \cite{hannan2025docslm}. Fully developed staged fusion is rare in the surveyed corpus.
% % ─────────────────────────────────────────────────────────────────────────────
% Auto-generated by summary_models_convert.py — do NOT edit manually.
% Re-generate:  python3 summary_models_convert.py models_summary.csv models_summary.tex
%
% Required packages (add to your preamble):
%   \usepackage{booktabs}
%   \usepackage{longtable}
%   \usepackage{array}
%   \usepackage{xcolor}
%   \usepackage{colortbl}
%   \usepackage{caption}
%   \usepackage{adjustbox}   % for \begin{adjustbox}
%   \usepackage{natbib}      % for \citeyear  (or use biblatex)
%   \usepackage{lscape}      % optional: landscape
% ─────────────────────────────────────────────────────────────────────────────

\definecolor{sectionbg}{gray}{0.88}

% ── tab:model_survey_ocr ──────────────────────────────────────────────────────────────
\begin{table*}[t]
\centering
% \small
\scriptsize
\captionsetup{font=small}
\setlength{\tabcolsep}{3pt}
\renewcommand{\arraystretch}{0.86}
\begin{adjustbox}{width=\textwidth}
\begin{tabular}{l l c l l l c c c c}
\toprule
\textbf{Model} & \textbf{Venue} & \textbf{OCR Note} & \textbf{Architecture} & \textbf{Vision Encoder} & \textbf{LLM Backbone} & \textbf{Mod.} & \textbf{Reasoning} & \textbf{Search} & \textbf{Agentic} \\
\midrule

% ── OCR-Free ─────────────────────────────────────────────────────────────
\rowcolor{sectionbg}
\multicolumn{10}{l}{\textbf{OCR-Free}} \\
\midrule
  {Docopilot}~(\citeyear{duan2025docopilot}) & CVPR & Train & MLLM & InternViT-300M & Multiple (Intern) & V & - & - & - \\
  {DocR1}~(\citeyear{xiong2026docr1}) & AAAI & -- & MLLM & Qwen2.5-VL & Qwen2.5-VL-7B & V & CoT, HD & - & - \\
  {DocSeeker}~(\citeyear{yan2025docseeker}) & CVPR & -- & MLLM & Qwen2.5-VL ViT & Qwen2.5-VL-7B-Instruct & V & CoT, HD & - & - \\
  {Leopard}~(\citeyear{jia2024leopard}) & TMLR & -- & MLLM & SigLIP-SO400M & LLaMa-3.1-8B, Mistral-7B & V & CoT & - & - \\
  {mPLUG-DocOwl2}~(\citeyear{hu2025mplug2}) & ACL & -- & MLLM & ViT-L/14 & mPLUG-Owl2 & V & - & - & - \\
  {Texthawk2}~(\citeyear{yu2024texthawk2}) & preprint & Train & MLLM & SigLIP-SO400M & Qwen2-7B-Instruct & V & - & - & - \\
  {URaG}~(\citeyear{shi2026urag}) & AAAI & -- & MLLM & Qwen2.5-VL ViT & Qwen2.5-VL-3B / Qwen2.5-VL-7B & V & HD & - & - \\
  {AVIR}~(\citeyear{li2025avir}) & MMAsia & -- & RA & Pix2Struct & Qwen2.5-VL-3B-AWQ & V & - & D & 1A \\
  {DFVC}~(\citeyear{qian2026accurate}) & MDPI & -- & RA & VisRAG, ColQwen & Qwen2-VL-2B-Instruct & V & - & D & - \\
  {DREAM}~(\citeyear{zhang2025dream}) & MM & -- & RA & InternVL2 ViT & InternVL2-4B & V & HD & J & 1A \\
  {M3DocRAG}~(\citeyear{cho2024m3docrag}) & preprint & -- & RA & ColPali, ColQwen & Multiple (Qwen/Intern/Idefics) & V & - & D & 1A \\
  {MM-R5}~(\citeyear{xu2026mmr5}) & AAAI & -- & RA & Qwen2.5-VL ViT & Qwen2.5-VL-7B & V & CoT & D & - \\
  {MoLoRAG}~(\citeyear{wu2025molorag}) & EMNLP & -- & RA & ColPali & Multiple (Qwen/DS/LLaVA) & V & - & D, Nav & - \\
  {Self-Attention Scoring}~(\citeyear{kang2024multi}) & ICDAR & -- & RA & Pix2Struct & Pix2Struct & V & - & D & 1A \\
  {SLEUTH}~(\citeyear{liu2025sleuth}) & preprint & -- & RA & ColPali, Qwen3-VL ViT & Multiple (Qwen/GLM/Gemini) & V & HD & D & 4A \\
  {SV-RAG}~(\citeyear{chen2025svrag}) & ICLR & -- & RA & - & Multiple (Intern/Pali/Phi) & V & - & D & 1A \\
  {VDocRAG}~(\citeyear{tanaka2025vdocrag}) & CVPR & Train & RA & Phi3V-4.2B & Phi3V-4.2B & V & - & D & 1A \\
  {Doc-V*}~(\citeyear{zheng2026docvstar}) & preprint & -- & AT & Qwen2.5-VL & Qwen2.5-VL-7B-Instruct & V & ReAct, HD & D, IR, QR, Nav & 1A \\
  {DocReact}~(\citeyear{wu2025docreact}) & ACL & -- & AT & ColPali, VisRAG & GPT-4o & V & ReAct, QD & D, IR, QR & 1A \\
  {MACT}~(\citeyear{yu2025mact}) & preprint & -- & AT & Multiple & Multiple (Qwen/Intern) & V & CoT, PE & D & 4A, SC, SV \\
  {VRAG-RL}~(\citeyear{wang2026vragrl}) & NeurIPS & -- & AT & Qwen2.5-VL ViT & Qwen2.5-VL-3B/7B-Instruct & V & ReAct, HD & D, IR, QR & 1A \\

\midrule[0.4pt]

% ── OCR-Used ─────────────────────────────────────────────────────────────
\rowcolor{sectionbg}
\multicolumn{10}{l}{\textbf{OCR-Used}} \\
\midrule
  {Arctic-TILT}~(\citeyear{borchmann2025arctic}) & ACL & -- & PDE & U-Net & T5-Large & T, V, L & - & - & - \\
  {GRAM}~(\citeyear{blau2024gram}) & CVPR & -- & PDE & DocFormerv2 & DocFormerv2 & T, V, L & - & - & - \\
  {Hi-VT5}~(\citeyear{tito2023hierarchical}) & PR & -- & PDE & DiT & T5-Base & T, V, L & - & - & - \\
  {RM-T5}~(\citeyear{dong2024multi}) & DAS & -- & PDE & DiT & T5-Base & T, V, L & - & - & - \\
  {CoR}~(\citeyear{guo2026end}) & preprint & Aux. & MLLM & Qwen2.5-VL & Qwen2.5-VL-7B & T, V & CoT, PE, HD & - & - \\
  {DocSLM}~(\citeyear{hannan2025docslm}) & CVPR & -- & MLLM & SigLIP2 & Qwen2.5-1.5B & T, V, L & - & - & - \\
  {InstructDoc}~(\citeyear{tanaka2024instructdoc}) & AAAI & -- & MLLM & CLIP & FlanT5, BLIP2 & T, V, L & - & - & - \\
  {LayTokenLLM}~(\citeyear{zhu2025laytokenllm}) & CVPR & -- & MLLM & - & LLaMA3-8B, Qwen1.5-7B & T, L & - & - & - \\
  {CREAM}~(\citeyear{zhang2024cream}) & ACM & -- & RA & Pix2Struct & LLaMA2-7B & T, V & HD & J & 1A \\
  {KGP}~(\citeyear{wang2024knowledge}) & AAAI & -- & RA & - & Multiple & T, L & - & S, IR, Nav & 1A \\
  {M2RAG}~(\citeyear{duan2025m2rag}) & NeurIPS & Aux. & RA & VisRAG-Ret & Multiple (Qwen) & T, V & - & J & 3A \\
  {MDocAgent}~(\citeyear{han2025mdocagent}) & preprint & Aux. & RA & ColPali, ColQwen2 & Multiple (Llama/Qwen) & T, V & - & J & 5A \\
  {MHier-RAG}~(\citeyear{gong2025mhier}) & preprint & -- & RA & Qwen-VL-Plus & Multiple (Qwen/DS/GPT) & T, V, L & CoT, HD & J & 1A \\
  {MLDocRAG}~(\citeyear{zhang2026mldocrag}) & preprint & -- & RA & - & Multiple (Qwen) & T, V, L & CoT, HD & D & 1A \\
  {MultiDocFusion}~(\citeyear{shin2025multidocfusion}) & ACL & -- & RA & - & Multiple (Qwen/LLaMa/Mistral) & T, V, L & HD & D & 1A \\
  {PDF-WuKong}~(\citeyear{xie2024wukong}) & preprint & Aux. & RA & IXC2-VL-4KHD & IXC2-VL-4KHD & T, V & - & J & 1A \\
  {RAG-DocVQA}~(\citeyear{lopez2025enhancing}) & ICDAR & Aux. & RA & DiT, Pix2Struct & Qwen2.5-VL-7B-Instruct & T, V, L & - & D & 1A \\
  {VisDoMRAG}~(\citeyear{suri2025visdom}) & NAACL & Aux. & RA & ColPali / ColQwen2 & Multiple(Qwen/Gemini/GPT) & T, V & CoT & J & 3A, SC \\
  {DMAP}~(\citeyear{fu2026dmap}) & WWW & Aux. & AT & ColPali & GPT-4o / Qwen-plus & T, V, L & HD & J & 2A, SR \\
  {DocAgent}~(\citeyear{sun2025docagent}) & EMNLP & Aux. & AT & -- & Multiple (GPT/Gemini/Claude) & T, V, L & ReAct, HD & J, IR, Nav & 3A \\
  {DocDancer}~(\citeyear{zhang2026docdancer}) & preprint & Aux. & AT & Qwen3-VL-235B & Qwen3-4B/30B-A3B & T, V, L & ReAct & J, IR, QR & 1A \\
  {DocLens}~(\citeyear{zhu2025doclens}) & preprint & -- & AT & -- & Multiple (GPT/Gemini/Claude) & T, V, L & CoT & J & 4A, SA \\
  {MM-Doc-R1}~(\citeyear{lin2026mmdocr1}) & preprint & Aux. & AT & Qwen2.5-VL ViT & Qwen3-4B/8B & T, V & ReAct, HD & S, IR, QR, Nav & 3A \\
  {SimpleDoc}~(\citeyear{jain2025simpledoc}) & EMNLP & Aux. & AT & ColQwen-2.5 & Multiple (Qwen) & T, V & - & J, IR, QR & 1A, SR, SV \\
  {ViDoRAG}~(\citeyear{wang2025vidorag}) & EMNLP & Aux. & AT & NV-Embed-V2, ColQwen2 & Multiple (Qwen/LLaMa/GPT) & T, V & ReAct, HD & J, IR & 3A, SR, SV \\

\bottomrule
\end{tabular}
\end{adjustbox}
\caption{Survey of MP-VRDU models grouped by OCR dependency. PDE\,=\,Page-to-Document Encoding Models; MLLM\,=\,MLLM-Centric Adaptation Models; RA\,=\,Retriever-Augmented Pipelines; AT\,=\,Adaptive-Trajectory Pipelines. T\,=\,text; V\,=\,visual; L\,=\,layout. CoT\,=\,Chain of Thought; ReAct\,=\,Reason\,+\,Act; PE\,=\,Plan-Execute; QD\,=\,Question Decomposition; HD\,=\,Hierarchical / Coarse-to-Fine. S\,=\,Sparse; D\,=\,Dense; J\,=\,Joint; IR\,=\,Iterative Retrieval; QR\,=\,Query Reformulation; Nav\,=\,Navigation. $n$A\,=\,$n$ agents; SR\,=\,Self-Refine; SV\,=\,Self-Verify; SC\,=\,Self-Consistency; SA\,=\,Sampling-Adjudication. OCR Note column: Aux.\,=\,model is not strictly OCR-dependent but incorporates OCR as an auxiliary framework component; Train\,=\,OCR-derived supervision is used during training but no OCR is used at inference. `--' indicates not reported / not applicable.}
\label{tab:model_survey_ocr}
\end{table*}

% \subsection{Model Survey by OCR Dependency}
% \label{app:ocr-survey}
% Table~\ref{tab:model_survey_ocr} reorganises the surveyed systems along the OCR-dependency axis to complement the architectural grouping in Table~\ref{tab:model_survey_arch}. The two views cut across each other rather than nesting: OCR-using systems span every architectural family, and OCR dependency is therefore largely orthogonal to architecture rather than a defining feature, as discussed in \S\ref{sec:architecture}. The ``OCR Note'' column flags hybrid cases in which OCR appears only as an auxiliary component or during training.
% Tables to prepare:
% 1. open source github table
% 2. resolution & adaptors/projectors table